% !TEX root = ../main.tex

\documentclass[../main.tex]{subfiles}

\graphicspath{{\subfix{../images/}}}

\begin{document}

В ходе выполнения курсовой работы разработана система визуализации
инерциальной навигации на основе полётного контроллера Pixhawk.
Реализованы следующие компоненты:
\begin{itemize}[leftmargin=2cm]
    \item Автономный модуль первичного считывания \texttt{RAW\_IMU} с записью
          данных в CSV для анализа и подбора калибровочных коэффициентов.
    \item Многопоточная архитектура C++14 с потокобезопасным буфером.
    \item Алгоритм калибровки сырых показаний с экспериментально
          определёнными смещениями нуля~\eqref{eq:biases}.
    \item Точное кватернионное интегрирование угловой скорости через
          ось-угол~\cite{kuipers1999} с нормализацией на каждом шаге.
    \item Вычисление матрицы вращения $\mathbf{R}\in SO(3)$ и углов Эйлера
          из кватерниона.
    \item Двойное интегрирование ускорения для оценки позиции~\cite{groves2013}.
    \item Механизм ZUPT~\cite{skog2010} для частичной компенсации дрейфа.
    \item OpenGL~3.3-визуализация: 3D-модель, координатная
          сетка, оси, траектория, HUD-панель
          с четырьмя графиками~\cite{glfw_docs}.
\end{itemize}
 
Алгоритм оценки ориентации показал высокое качество: 3D-модель точно
воспроизводила реальные повороты устройства. Алгоритм оценки позиции
продемонстрировал принципиальное ограничение МЭМС-датчиков~\cite{titterton2004}:
ошибка нарастает достаточно быстро~\eqref{eq:drift}, и получение
практически пригодных координат требует слияния данных IMU с корректирующими
устройствами, датчиками с применением фильтра Калмана~\cite{groves2013}.

\end{document}